% from ega2-1-fr.tex
\oldpage[II]{119}

\begin{proposition}[6.3.7]
\label{II.6.3.7-fr}
Chacun des $X'_{ij}$ est un $X$-préschéma intègre et normal, et son corps des
fonctions rationnelles $R(X'_{ij})$ s'identifie canoniquement à la fermeture
algébrique $K'_{ij}$ de $k(\xi_i)$ dans $K_{ij}$.
\end{proposition}

En vertu de ce qui précède, on peut supposer $X$ intègre, donc $r=1$, et
$s_1=1$, de sorte que l'unique algèbre $A_1$ est un corps $K$~; soit $\xi$ le
point générique de $X$, et soit $f:X'\to X$ le morphisme structural. Pour tout
ouvert affine non vide $U$ de $X$, $f^{-1}(U)$ s'identifie au spectre de la
fermeture intégrale $B'_U$ dans le corps $K$ de l'anneau intègre
$B_U=\Gamma(U,\mathcal{O}_X)$ (\hyperref[II.6.3.4-fr]{6.3.4})~; comme l'anneau
$B'_U$ est intègre et intégralement clos, il en est de même des anneaux locaux
des points de son spectre, donc $f^{-1}(U)$ est par définition un schéma
\emph{intègre et normal} ((\hyperref[0.4.1.4-fr]{0, 4.1.4}) et
(\hyperref[I.5.1.4-fr]{I, 5.1.4})). En outre, comme $(0)$ est le seul idéal
premier de $B'_U$ au-dessus de l'idéal premier $(0)$ de $B_U$
(\cite[t.~I, p.~259]{I-13}), $f^{-1}(\xi)$ se réduit à un seul point $\xi'$,
et $k(\xi')$ est le corps des fractions $K'$ de $B'_U$, qui n'est autre que la
fermeture algébrique de $k(\xi)$ dans $K$. Enfin, $X'$ est irréductible, car
lorsque $U$ parcourt l'ensemble des ouverts affines non vides de $X$, les
$f^{-1}(U)$ constituent un recouvrement ouvert de $X'$ formé d'ouverts
irréductibles~; en outre l'intersection $f^{-1}(U\cap V)$ de deux de ces
ouverts contient $\xi'$, donc n'est pas vide, et on conclut par
(\hyperref[0.2.1.4-fr]{0, 2.1.4}).

\begin{corollary}[6.3.8]
\label{II.6.3.8-fr}
Soit $X$ un préschéma réduit n'ayant qu'un nombre fini de composantes
irréductibles $X_i$ ($1\leq i\leq r$), et soit $\xi_i$ le point générique de
$X_i$. La fermeture intégrale $X'$ de $X$ relativement à $\mathcal{R}(X)$ est
somme de $r$ $X$-schémas $X'_i$ qui sont intègres et normaux. Si $f:X'\to X$
est le morphisme structural, $f^{-1}(\xi_i)$ se réduit au point générique
$\xi'_i$ de $X'_i$ et on a $k(\xi'_i)=k(\xi_i)$, autrement dit $f$ est
birationnel.
\end{corollary}

On dit dans ce cas que $X'$ est le \emph{normalisé} du préschéma réduit $X$~;
on notera que $f$, étant birationnel et entier, est \emph{surjectif}
(\hyperref[II.6.1.10-fr]{6.1.10}). Pour que l'on ait alors $X'=X$, il faut et
il suffit que $X$ soit normal. Lorsque $X$ est un préschéma intègre, il résulte
de (\hyperref[II.6.3.8-fr]{6.3.8}) que son normalisé $X'$ est intègre.

\begin{env}[6.3.9]
\label{II.6.3.9-fr}
Soient $X$, $Y$ deux préschémas intègres, $f:X\to Y$ un morphisme dominant,
$L=R(X)$, $K=R(Y)$ les corps des fonctions rationnelles de $X$ et $Y$~; il
correspond canoniquement à $f$ une injection $K\to L$, et lorsqu'on identifie
$K$ (resp. $L$) au faisceau simple $\mathcal{R}(Y)$ (resp.
$\mathcal{R}(X)$), cette injection est un $f$-morphisme. Soit $K_1$ (resp.
$L_1$) une extension de $K$ (resp. $L$) et supposons donné un monomorphisme
$K_1\to L_1$ tel que le diagramme
\[
\xymatrix{
K_1\ar[r]&L_1\\
K\ar[u]\ar[r]&L\ar[u]
}
\]
soit commutatif~; si $K_1$ (resp. $L_1$) est considéré comme faisceau simple
sur $Y$ (resp. $X$), donc comme une $\mathcal{R}(Y)$-Algèbre (resp. une
$\mathcal{R}(X)$-Algèbre), cela signifie que $K_1\to L_1$ est un
$f$-morphisme. Cela étant, si $X'$ (resp. $Y'$) est la fermeture intégrale de
$X$ (resp. $Y$) relativement à $L_1$ (resp. $K_1$), $X'$ (resp. $Y'$) est un
préschéma intègre et normal (\hyperref[II.6.3.6-fr]{6.3.6}) dont le corps des
fonctions rationnelles s'identifie canoniquement à la fermeture algébrique

